\documentclass[12pt]{article}

\usepackage[utf8]{inputenc}
\usepackage[T1]{fontenc}
\usepackage{amsmath, amssymb, amsthm}
\usepackage{geometry}
\geometry{a4paper, margin=2.5cm}

\title{Teorema lui Weierstrass}
\author{}
\date{}

\theoremstyle{plain}
\newtheorem{theorem}{Teoremă}

\begin{document}

\maketitle

\begin{theorem}
\textbf{Teorema lui Weierstrass.}  
Fie $\Omega$ o mulțime deschisă din $\mathbb{C}$ și fie $(f_n)$ un șir de funcții olomorfe în $\Omega$.  
Presupunem că $(f_n)$ converge uniform pe orice submulțime compactă $K \subset \Omega$ către o funcție $f$.

Atunci funcția $f$ este olomorfă în $\Omega$, iar pentru orice $p \ge 0$ are loc



\[
f_n^{(p)} \to f^{(p)} \quad \text{uniform pe orice compact } K \subset \Omega.
\]



Cu alte cuvinte, convergența uniformă pe compacte a funcțiilor olomorfe implică convergența uniformă pe compacte a tuturor derivatelor lor.
\end{theorem}

\end{document}
